\section{Lorentzian polynomials and their apolar algebra}

We collect useful facts on Lorentzian polynomials in order to simplify the analysis in the next section.

Let $R=\R[X_1,\ldots,X_n]$ and $S=\R[x_1,\ldots,x_n]$ be two polynomial rings. Following the idea of Macaulay's \textit{inverse systems}, we define an action of $R$ on $S$ by
\[
    f(X_1,\ldots,X_n)\cdot Q\coloneqq  f(\frac{\partial}{\partial x_1},\ldots,\frac{\partial}{\partial x_n})Q,
\]
and we use $\ann(Q)$ to denote the ideal
\[
    \{f\in R\mid f\cdot Q=0\}. 
\]
When $0\neq Q\in S$ is homogeneous of degree $d\geq 2$, we define its \textit{apolar algebra} as the graded ring
\[
    A^\bullet_Q\coloneqq R/\ann(Q),
\]
which has Poincar\'e duality with $\deg\colon A^d_Q\to \R$ given by $\deg(f)=f\cdot Q$.

Let $(A^\bullet,\deg)$ be a Poincar\'e duality algebra and elements $\chi_1,\ldots,\chi_n$ be elements in $A^1$. We define the \textit{volume polynomial} with respect to $\chi_1,\ldots,\chi_n$ by
\begin{align*}
    Q(x_1,\ldots,x_n)=&\frac{1}{d!}\deg\big((x_1\chi_1+\cdots+x_n\chi_n)^d\big)\\
    =&\sum_{\substack{
    k_1+\cdots +k_n=d\\k_1,\ldots ,k_n\geq 0    }}\deg(\chi_1^{k_1}\cdots\chi_n^{k_n})\frac{x_1^{k_1}\cdots x_n^{k_n}}{k_1!\cdots k_n!}.
\end{align*}
We have the following correspondence.
\begin{theorem}[see for example, {\cite[Theorem 21.6 and Exercise 21.7]{Eis95}}] The map
\[
    \big\{0\neq Q\in S^d\big\}\xrightarrow{Q\mapsto A^\bullet_Q}\big\{\substack{\text{graded quotient rings of}\ R\ \text{with a degree map in degree}\ d\\ \text{satisfying Poincar\'e duality}}\big\},
\]
is bijective, with the inverse given by taking the volume polynomial with respect to the images of $X_1,\ldots,X_n$.
\end{theorem}

A homogeneous polynomial of degree $d$ is \textit{strictly Lorentzian} if it has positive coefficients and the Hessian matrices of its $(d-2)$-th order partial derivatives have the Lorentzian signature; a polynomial is \textit{Lorentzian} if it's the limit of strictly Lorentzian polynomials. For a homogeneous polynomial $Q\neq 0$, we define $K_Q\coloneqq \{a_1X_1+\cdots+a_nX_n\mid a_1,\ldots,a_n>0\}\subseteq A^1_Q$. We say a Poincar\'e duality algebra satisfies \textit{mixed $(\mathrm{HR}^{\leq m})$} with respect to a cone $K$ if the mixed Hodge-Riemann relations are satisfied for $0\leq k\leq m$. We have the following characterization of Lorentzian polynomials:

\begin{theorem}[{\cite[Theorem 3.2]{Eur21}}]\label{thm: lorentzian and HR}
A nonzero homogeneous polynomial $Q\in S$  is Lorentzian if and only if the triple $(A^\bullet_Q,\deg,K_Q)$ satisfies mixed $(\mathrm{HR}^{\leq 1})$.
\end{theorem}

Let $d\Delta_n=\{k=(k_1,\ldots,k_n)\in \mathbb Z_{\geq 0}^n\mid \sum_{i=1}^n k_i=d\}$. For $Q\in S^d$, the support of $Q$, denoted $\supp Q$, is the set of exponent vectors where the corresponding coefficient of $Q$ is nonzero. A subset $\mathcal S$ of $d\Delta_n$ is \textit{M-convex} if for any $k,k'\in \mathcal S$ and for any $i$ such that $k_i<k'_i$, there exists $j$ such that $k_j>k'_j$ and both
\[
    k+e_i-e_j\text{ and } k'-e_i+e_j
\]
are contained in $\mathcal S$. We also need the following characterization of Lorentzian polynomials.

\begin{theorem}[{\cite[Theorem 2.25]{BH20}}]\label{thm: Lorentzian-criterion}
A homogeneous polynomial $Q\in S^d$ is Lorentzian if and only if
\begin{enumerate}
    \item  Q has nonnegative coefficients,
    \item $\supp(Q)\subseteq d\Delta_n$ is M-convex, and
    \item For any $k_1+\cdots +k_n=d-2,\ k_1,\ldots,k_n\geq 0$, the Hessian matrix of \[\left(\frac{\partial}{\partial x_1}\right)^{k_1}\cdots\left(\frac{\partial}{\partial x_n}\right)^{k_n}Q \] has at most one positive eigenvalue.
\end{enumerate}
\end{theorem}

We say a polynomial $Q\in S$ is \textit{concise} if $R^1\cong A^1_Q$, i.e., $f\cdot Q\neq 0$ for any nonzero linear form $f\in R$. The following variation of \cite[Theorem 2.10]{BH20} would be useful for proving the strict Lorentzian property.

\begin{proposition}\label{prop: strictification}
   Let $(A^\bullet,\deg,K)$ be a Poincar\'e duality algebra satisfying \textit{mixed $(\mathrm{HR}^{\leq 1})$} with respect to $K$.
   \begin{enumerate}
   \item If $\gamma_1,\ldots,\gamma_m\in\overline{K}$, then their volume polynomial is Lorentzian;
   \item If $\gamma_1,\ldots,\gamma_m\in K$ are linearly independent, then their volume polynomial is strictly Lorentzian.
   \end{enumerate}
   In particular, for any rank $m$ real $n\times m$ matrix $M$ with positive entries and any concise Lorentzian polynomial $Q$ in $n$ variables, $Q(My)$ is a strict Lorentzian in $m$ variables.
\end{proposition}
\begin{proof}
    Consider the case when $\gamma_1,\ldots,\gamma_m\in K$. Let $F=F(y_1,\ldots,y_m)$ be the volume polynomial of $\gamma_1,\ldots,\gamma_m$. Mixed Hodge-Riemann relation in degree 0 implies
\[
    \deg\!\left(\prod_{i=1}^m\gamma_i^{k_i}\right)>0
\]
for any $k_1+\cdots+k_m=d,\ k_1,\ldots,k_m\geq 0$. In other words, $F$ has positive coefficients.

In degree 1, the Hodge-Riemann relation implies that the symmetric form
\[
    A^1\times A^1 \to \R, \ \ (\alpha,\beta)\mapsto \deg\left(\prod_{i=1}^m\gamma_i^{k_i}\cdot \alpha\beta\right) 
\]
has the Lorentzian signature for any $k_1+\cdots +k_m=d-2,\ k_1,\ldots,k_m\geq 0$. Cauchy's interlacing formula, plus the fact that this form is positive on $\gamma_i$, implies that the restriction of this form to the subspace $\operatorname{span}\{ \gamma_i\}_i$ also has the Lorentzian signature. By pulling back this form to $\R^m$ via $e_i\mapsto\gamma_i$, we obtain the Hessian matrix of $\left(\frac{\partial}{\partial y_1}\right)^{k_1}\cdots\left(\frac{\partial}{\partial y_m}\right)^{k_m}F$, so this matrix has exactly one positive eigenvalue. This proves that $F$ is Lorentzian, by \ref{thm: Lorentzian-criterion}.
When $\gamma_1,\cdots,\gamma_m\in \overline{K}$, a limiting argument shows that $F$ is still Lorentzian; when $\gamma_1,\ldots,\gamma_m\in K$ are linearly independent, one further sees that the Hessian matrices have the Lorentzian signature and therefore $F$ is strictly Lorentzian.
    Given a concise Lorentzian polynomial $Q$, consider its apolar algebra with the standard generators $\chi_1,\ldots,\chi_n$. Let $\gamma_1,\ldots,\gamma_m$ be elements in $A^1$ defined by $\gamma=M^T\chi$. One sees that 
    \[
    \deg(f(\gamma))=\deg(f(M^T\chi))=f(M^TX)\cdot Q(x)=f(Y)\cdot Q(My)
    \]
    for any homogeneous polynomial $f$ of degree $d$, where $X_i$ and $Y_j$ act as partial differentials. Since $M$ has positive entries and that $M$ is of rank $m$, $\gamma$'s are linearly independent elements in $K$ and $Q(My)$ is precisely their volume polynomial. Therefore, $Q(My)$ is strictly Lorentzian by the previous claim.



\end{proof}

\begin{remark}
    The proposition above also shows that $A^\bullet_{Q(My)}$ satisfies mixed Hodge-Riemann relations in degree less than or equal to $1$ for classes in $\overline K_{Q(My)}\setminus 0$, by \cref{thm: lorentzian and HR}. This implication cannot extend to higher Hodge-Riemann relations. See \cite[Example 6.17]{NW26}, for example.
\end{remark}

